% C087 | 双曲线焦点弦长公式与取值范围 | 简短源稿（非独立文档）
\textbf{C087\quad 双曲线焦点弦长公式与取值范围}

设 $H:\dfrac{x^2}{a^2}-\dfrac{y^2}{b^2}=1$（$a,b>0$），
$c^2=a^2+b^2$，直线 $AB$ 经过焦点 $F_2(c,0)$ 且与 $H$ 有两个不同有限交点。
当直线斜率为 $k$ 时，必须 $k\ne\pm b/a$，此时
\[
\boxed{|AB|=\frac{2ab^2(1+k^2)}{|a^2k^2-b^2|}}.
\]
以倾斜角 $\theta\in[0,\pi)$ 表示则可包含竖直弦：
\[
\boxed{|AB|=\frac{2ab^2}{|a^2-c^2\cos^2\theta|}},
\qquad \cos\theta\ne\pm\frac ac.
\]
\textbf{异支弦：}$|k|<b/a$，$|AB|\in[2a,+\infty)$，
在实轴 $y=0$ 取最小值 $2a$。
\textbf{同支弦：}$|k|>b/a$ 或直线竖直，
$|AB|\in[2b^2/a,+\infty)$，在 $x=c$ 取最小值 $2b^2/a$。
\textbf{趋近边界：}当 $k\to\pm b/a$，弦长趋于无穷大；
$k=\pm b/a$ 时直线平行渐近线，只剩一个有限交点，不能称作焦点弦。
\textbf{速证：}令 $t=x-c$，代入 $y=kt$ 得
$(b^2-a^2k^2)t^2+2cb^2t+b^4=0$，判别式
$\Delta=4a^2b^4(1+k^2)$，于是 $|AB|=\sqrt{1+k^2}\sqrt\Delta/|b^2-a^2k^2|$。
